\documentclass[10pt]{beamer}

% --- Modern Theme ---
\usetheme{metropolis}
\usecolortheme{default}
\setbeamercolor{background canvas}{bg=white}
\metroset{sectionpage=none,subsectionpage=none}

% --- Packages ---
\usepackage[utf8]{inputenc}
\usepackage{booktabs}
\usepackage{array}
\usepackage{tabularx}
\usepackage{amsmath}
\usepackage{amssymb}
\usepackage[scale=2]{ccicons}
\usepackage{pgfplots}
\usepgfplotslibrary{dateplot}
\usepackage{xspace}

% --- Title Information ---
\title{Foundations of Artificial Intelligence}
\subtitle{Part III: Knowledge Representation \& Inference}
\institute{Department of Computer Science}
\date{}

\begin{document}

\maketitle

\begin{frame}{Lecture Roadmap}
    \setbeamertemplate{section in toc}[sections numbered]
    \tableofcontents[hideallsubsections]
\end{frame}

\begin{frame}{Learning Objectives}
    By the end of this lecture, you should be able to:
    \begin{itemize}
        \item Explain why representation choices determine what an AI system can infer.
        \item Model facts and rules in propositional and first-order logic.
        \item Apply inference ideas: entailment, chaining, and resolution at a conceptual level.
        \item Relate rule systems, semantic networks, and ontologies to real AI applications.
        \item Analyze the strengths/limits of each representation for practical tasks.
    \end{itemize}
\end{frame}

\section{Why Knowledge Representation?}

\begin{frame}{From Data to Knowledge}
    AI systems need more than raw data:
    \begin{itemize}
        \item \textbf{Data:} isolated observations (e.g., temperature is $39^{\circ}\mathrm{C}$).
        \item \textbf{Information:} structured observations (patient has fever + cough).
        \item \textbf{Knowledge:} reusable, inferable structure (if fever and cough, suspect infection).
    \end{itemize}
    \vspace{0.2cm}
    \textit{Book perspective: a representation is useful only if it supports effective inference.}
\end{frame}

\begin{frame}{Real-World Motivation}
    \begin{itemize}
        \item \textbf{Clinical decision support:} map symptoms to diagnostic hypotheses.
        \item \textbf{Fraud detection:} infer suspicious patterns from transaction relations.
        \item \textbf{Customer support bots:} apply business rules consistently.
        \item \textbf{Industrial maintenance:} reason over component dependencies and failure rules.
    \end{itemize}
    \vspace{0.25cm}
    Common requirement: \textbf{explicit knowledge} + \textbf{traceable reasoning}.
\end{frame}

\begin{frame}{What Makes a Good Representation?}
    We judge KR by four practical criteria:
    \begin{enumerate}
        \item \textbf{Expressiveness:} Can it describe what matters?
        \item \textbf{Inference support:} Can we derive useful conclusions?
        \item \textbf{Computational cost:} Can we reason in acceptable time?
        \item \textbf{Maintainability:} Can humans update and validate it?
    \end{enumerate}
\end{frame}

\section{Logic Foundations}

\begin{frame}{Propositional Logic: Core Idea}
    Propositional logic is the \textbf{simplest formal language} for expressing and reasoning about truth.
    \begin{itemize}
        \item \textbf{Atomic propositions} are the basic building blocks — indivisible statements that are either \emph{true} or \emph{false} (e.g., $P$ = "It is raining", $Q$ = "The road is wet").
        \item \textbf{Connectives} combine propositions into complex formulas:
            $\neg$ (not), $\land$ (and), $\lor$ (or), $\Rightarrow$ (if-then), $\Leftrightarrow$ (if and only if).
        \item A \textbf{model} (or interpretation) assigns True/False to every atom; a formula is evaluated under this assignment.
        \item \textbf{Limitation:} Propositions are \emph{flat} — no objects, no variables, no structure inside. Every distinct fact needs its own symbol.
    \end{itemize}
    \vspace{0.15cm}
    \textbf{Smart-home example:}
    \[
        (\text{Smoke} \land \text{Heat}) \Rightarrow \text{Alarm}
    \]
    If both sensors trigger, sound the alarm. This single rule replaces ad-hoc 
    if/else code and is \emph{inspectable, explainable, and reusable}.
\end{frame}

\begin{frame}{Semantics (1/2): Core Concepts}
    \textbf{Semantics} defines \emph{what a formula means} — its truth value across all possible worlds (models).

    \vspace{0.2cm}
    \begin{itemize}
        \item A \textbf{model} (or interpretation) assigns \emph{True} or \emph{False} to every atomic proposition.
              Think of it as one possible ``state of the world.''
        \item If $KB$ uses $n$ atoms, there are $2^n$ possible models to consider.
              With 3 atoms, that is 8 worlds; with 10 atoms, 1024 worlds.
        \item A formula is \textbf{true under a model} if the assignment makes it evaluate to True.
        \item A model \textbf{satisfies} $KB$ if \emph{every} sentence in $KB$ is true under it.
    \end{itemize}

    \vspace{0.3cm}
    \textit{Semantics is what separates logic from mere symbol manipulation — it grounds every symbol in meaning.}
\end{frame}

\begin{frame}{Semantics (2/2): Truth Table Example}
    \textbf{Example:} Let $P$ = ``it is raining'' and $Q$ = ``the road is wet.''

    \vspace{0.2cm}
    We evaluate the formula $P \Rightarrow Q$ across all four possible models:

    \vspace{0.15cm}
    \begin{center}
    \renewcommand{\arraystretch}{1.3}
    \begin{tabular}{cc|cl}
        \toprule
        $P$ & $Q$ & $P \Rightarrow Q$ & Interpretation \\
        \midrule
        T & T & T & Rained and road is wet — rule holds. \\
        T & F & \alert{F} & Rained but road is dry — \alert{rule violated!} \\
        F & T & T & No rain, road wet (other cause) — rule holds. \\
        F & F & T & No rain, road dry — rule holds vacuously. \\
        \bottomrule
    \end{tabular}
    \end{center}

    \vspace{0.2cm}
    The implication fails in \textbf{only one} model — when the antecedent is true but the consequent is false.

    \vspace{0.1cm}
    \textit{A formula is semantically meaningful precisely because we can identify which worlds satisfy it.}
\end{frame}

\begin{frame}{Entailment}
    \small
    \textbf{Entailment} ($KB \models \alpha$) is the central concept of logical reasoning:
    \[
        KB \models \alpha \quad \iff \quad \text{every model satisfying } KB \text{ also satisfies } \alpha
    \]
    Intuitively: \emph{``If everything in $KB$ is true, then $\alpha$ must be true too — no escape.''}

    \vspace{0.08cm}
    \textbf{Example:}
    \begin{itemize}\setlength{\itemsep}{1pt}
        \item $KB = \{ P \Rightarrow Q,\; P \}$
        \item In every model where both $P$ and $P \Rightarrow Q$ hold, $Q$ must hold.
        \item Therefore: $KB \models Q$.
    \end{itemize}

    \vspace{0.08cm}
    \textbf{Entailment vs.\ Derivation — a critical distinction:}
    \begin{itemize}\setlength{\itemsep}{2pt}
        \item $KB \models \alpha$ \;(entailment) — a \emph{semantic} claim: about truth across all worlds.
        \item $KB \vdash \alpha$ \;(derivation) — a \emph{syntactic} claim: produced by applying inference rules.
        \item \textbf{Sound} system: $KB \vdash \alpha \;\Rightarrow\; KB \models \alpha$ \quad (never derives falsehoods)
        \item \textbf{Complete} system: $KB \models \alpha \;\Rightarrow\; KB \vdash \alpha$ \quad (misses nothing true)
    \end{itemize}

    \vspace{0.06cm}
    \textit{Entailment is the gold standard. Inference procedures try to match it efficiently.}
\end{frame}

\begin{frame}{Validity}
    \small
    A sentence is \textbf{valid} (a \emph{tautology}) if it is true in \textbf{every possible model} — regardless of what the atoms mean.

    \vspace{0.06cm}
    \textbf{Examples of valid sentences:}
    \begin{itemize}\setlength{\itemsep}{1pt}
        \item $P \lor \neg P$ \;— always true (law of excluded middle).
        \item $(P \Rightarrow Q) \Leftrightarrow (\neg Q \Rightarrow \neg P)$ \;— contrapositive equivalence.
        \item $\neg (P \land \neg P)$ \;— contradiction can never hold (law of non-contradiction).
    \end{itemize}

    \vspace{0.06cm}
    \textbf{Three-way classification of sentences:}
    \begin{center}
    \renewcommand{\arraystretch}{1.1}
    \begin{tabular}{>{\bfseries}lll}
        \toprule
        Type & True in & Example \\
        \midrule
        Valid (tautology) & Every model & $P \lor \neg P$ \\
        Satisfiable & At least one model & $P \land Q$ \\
        Unsatisfiable & No model & $P \land \neg P$ \\
        \bottomrule
    \end{tabular}
    \end{center}

    \vspace{0.06cm}
    \textbf{Why validity matters for inference:}\\
    Resolution provers prove $KB \models \alpha$ by showing $KB \land \neg\alpha$ is \emph{unsatisfiable} —
    turning an entailment question into a satisfiability check.
\end{frame}

\begin{frame}{Mini Propositional Example}
    \textbf{Scenario:} A delivery system has two rules about road conditions.

    \vspace{0.1cm}
    \textbf{Knowledge Base (KB):}
    \[
    \begin{aligned}
      &\text{(1)}\quad R \Rightarrow W \quad \text{(``If it rained, the road is wet.'')} \\
      &\text{(2)}\quad W \Rightarrow D \quad \text{(``If the road is wet, deliveries are delayed.'')} \\
      &\text{(3)}\quad R \quad\quad\quad \text{(``It rained.'' — observed fact)}
    \end{aligned}
    \]

    \vspace{0.1cm}
    \textbf{Query:} Is $D$ (delivery delayed) entailed?

    \vspace{0.1cm}
    \textbf{Step-by-step inference via Modus Ponens:}
    \[
      \frac{R \Rightarrow W,\quad R}{W} \qquad \text{then} \qquad \frac{W \Rightarrow D,\quad W}{D}
    \]

    \textbf{Result:} $KB \vdash D$, and since the system is sound, $KB \models D$.

    \vspace{0.1cm}
    \textit{Insight: Modus Ponens fires rules like a chain. Each step produces a new fact that can trigger the next rule. This is \textbf{forward chaining} in action.}
\end{frame}

\begin{frame}{CNF and SAT (1/2): Core Idea}
        \textbf{CNF (Conjunctive Normal Form)} is a standard way to write formulas.

        \vspace{0.1cm}
        \begin{itemize}
                \item \textbf{Literal:} $A$ or $\neg A$
                \item \textbf{Clause:} OR of literals, e.g., $(\neg A \lor B)$
                \item \textbf{CNF:} AND of clauses
        \end{itemize}

        \vspace{0.1cm}
        Example CNF:
        \[
            (\neg A \lor B) \land (\neg B \lor C) \land A
        \]

        \vspace{0.1cm}
        \textbf{Why CNF?}
        \begin{itemize}
                \item SAT solvers are optimized for this shape.
                \item Uniform format makes reasoning faster in practice.
        \end{itemize}
\end{frame}

\begin{frame}{CNF and SAT (2/2): SAT in Practice}
        \textbf{SAT question:} Is there \emph{any} assignment that makes the CNF true?

        \vspace{0.12cm}
        \[
            (\neg A \lor B) \land (\neg B \lor C) \land A
        \]

        \vspace{0.08cm}
        Try $A{=}T, B{=}T, C{=}T$:
        \begin{itemize}
                \item $(\neg A \lor B) = (F \lor T) = T$
                \item $(\neg B \lor C) = (F \lor T) = T$
                \item $A = T$
        \end{itemize}
        So this formula is \textbf{satisfiable}.

        \vspace{0.1cm}
        \textbf{Where SAT is used:}
        \begin{itemize}
                \item planning,
                \item hardware/software verification,
                \item scheduling and timetabling.
        \end{itemize}
\end{frame}

\begin{frame}{First-Order Logic (FOL) (1/2): Why We Need It}
        \textbf{Problem with propositional logic:}
        \begin{itemize}
                \item It cannot naturally talk about \emph{objects} and \emph{relations}.
                \item You need one symbol per fact (not scalable).
        \end{itemize}

        \vspace{0.12cm}
        Example idea: ``Every patient with fever needs rest."
        \begin{itemize}
                \item In propositional logic: many separate symbols/rules.
                \item In FOL: one reusable rule for all patients.
        \end{itemize}

        \vspace{0.12cm}
        \textbf{Goal of FOL:} write compact rules that generalize to many entities.
\end{frame}

\begin{frame}{First-Order Logic (FOL) (2/2): Building Blocks}
        FOL adds structure using these elements:
        \begin{itemize}
                \item \textbf{Constants:} Ali, London, DrSmith
                \item \textbf{Variables:} $x, y, z$
                \item \textbf{Predicates:} $Doctor(x)$, $Treats(x,y)$
                \item \textbf{Quantifiers:} $\forall$ (for all), $\exists$ (there exists)
        \end{itemize}

        \vspace{0.12cm}
        Example rule:
        \[
            \forall p\; (Patient(p) \land HasFever(p) \Rightarrow NeedsRest(p))
        \]

        \vspace{0.08cm}
        One sentence applies to every patient.
\end{frame}

\begin{frame}{FOL Example (1/2): Book-Style}
        \textbf{Classic example:}
        \[
            \forall x\; (Human(x) \Rightarrow Mortal(x))
        \]
        \[
            Human(Socrates)
        \]

        \vspace{0.1cm}
        Therefore:
        \[
            Mortal(Socrates)
        \]

        \vspace{0.1cm}
        This is: apply a general rule to one specific object.
\end{frame}

\begin{frame}{FOL Example (2/2): Real-World}
        \textbf{Healthcare rule:}
        \[
            \forall p\; (HasFlu(p) \Rightarrow NeedsRest(p))
        \]
        \[
            HasFlu(Ali)
        \]

        \vspace{0.1cm}
        Therefore:
        \[
            NeedsRest(Ali)
        \]

        \vspace{0.12cm}
        \textbf{Takeaway:} same pattern as Socrates, but now in a practical domain.
\end{frame}

\begin{frame}{Unification (1/2): Intuition}
        \textbf{Unification} means:
        find substitutions that make two expressions match.

        \vspace{0.12cm}
        Example:
        \[
            Treats(DrSmith, x) \quad \text{and} \quad Treats(DrSmith, Sam)
        \]
        Unifier:
        \[
            \theta = \{x/Sam\}
        \]

        \vspace{0.12cm}
        After substitution, both expressions are identical.
\end{frame}

\begin{frame}{Unification (2/2): Success and Failure}
        \textbf{Success:}
        \[
            Knows(x,y),\; Knows(Alice,Bob)
            \Rightarrow \theta = \{x/Alice,\; y/Bob\}
        \]

        \vspace{0.1cm}
        \textbf{Failure:}
        \[
            Treats(DrSmith,Sam),\; Treats(DrJones,Sam)
        \]
        cannot unify because $DrSmith \neq DrJones$.

        \vspace{0.1cm}
        \textbf{Why important:}
        \begin{itemize}
                \item backward chaining,
                \item resolution,
                \item Prolog-style rule matching.
        \end{itemize}
\end{frame}

\section{Inference in Knowledge Systems}

\begin{frame}{Inference: Big Picture}
    Inference answers: \textit{What new facts follow from what we already know?}

    \vspace{0.2cm}
    Desirable properties:
    \begin{itemize}
        \item \textbf{Soundness:} never derive false conclusions.
        \item \textbf{Completeness:} derive every true entailed conclusion (within logic limits).
        \item \textbf{Efficiency:} finish in practical time.
    \end{itemize}
\end{frame}

\begin{frame}{Forward vs Backward Chaining}
    \begin{columns}[T,onlytextwidth]
        \column{0.48\textwidth}
        \alert{Forward Chaining}
        \begin{itemize}
            \item Data-driven
            \item Fire rules from known facts
            \item Good for streaming/event systems
        \end{itemize}

        \column{0.48\textwidth}
        \alert{Backward Chaining}
        \begin{itemize}
            \item Goal-driven
            \item Start from query, prove prerequisites
            \item Good for diagnosis/query answering
        \end{itemize}
    \end{columns}
\end{frame}

\begin{frame}{Resolution (Conceptual)}
    Resolution works by contradiction:
    \begin{enumerate}
        \item Convert $KB$ and $\neg Query$ into CNF clauses.
        \item Repeatedly resolve complementary literals.
        \item Derive empty clause to prove entailment.
    \end{enumerate}
    \vspace{0.2cm}
    Why useful:
    \begin{itemize}
        \item Uniform rule for automated theorem proving.
        \item Foundation for many symbolic reasoners.
    \end{itemize}
\end{frame}

\section{Rules, Semantic Networks, Ontologies}

\begin{frame}{Production Rules}
    Rules encode procedural domain knowledge:
    \[
      \text{IF condition THEN conclusion/action}
    \]
    \vspace{0.2cm}
    Example (e-commerce):
    \begin{itemize}
        \item IF customer is premium AND cart value $>$ 500 THEN shipping = free.
        \item IF payment fails 3 times THEN flag transaction.
    \end{itemize}
\end{frame}

\begin{frame}{Rule-Based System Architecture}
    \begin{itemize}
        \item \textbf{Knowledge Base:} facts + rules.
        \item \textbf{Working Memory:} current case-specific facts.
        \item \textbf{Inference Engine:} chooses and applies rules.
        \item \textbf{Explanation Layer:} why a decision was reached.
    \end{itemize}
    \vspace{0.2cm}
    \textit{Strength: interpretable decisions. Limitation: rule maintenance can grow hard at scale.}
\end{frame}

\begin{frame}{Semantic Networks (Graph View of Knowledge)}
    \begin{itemize}
        \item Nodes represent concepts/instances.
        \item Edges represent relations: \texttt{is-a}, \texttt{part-of}, \texttt{causes}, etc.
        \item Supports inheritance-like reasoning.
    \end{itemize}
    \vspace{0.2cm}
    Example:
    \begin{itemize}
        \item \texttt{Bird is-a Animal}
        \item \texttt{Canary is-a Bird}
        \item infer: \texttt{Canary is-a Animal}
    \end{itemize}
\end{frame}

\begin{frame}{Logic vs Rules vs Semantic Nets}
    \footnotesize
    \setlength{\tabcolsep}{4pt}
    \renewcommand{\arraystretch}{1.13}
    \begin{tabularx}{\textwidth}{>{\raggedright\arraybackslash}p{2.5cm} *{3}{>{\raggedright\arraybackslash}X}}
        \toprule
        \textbf{Approach} & \textbf{Main Strength} & \textbf{Typical Use} & \textbf{Main Limitation} \\
        \midrule
        Logic (PL/FOL) & Precise semantics & Formal reasoning, verification & Can become computationally heavy \\
        Production Rules & Simple expert knowledge encoding & Policy, diagnosis, decision support & Rule explosion/conflicts \\
        Semantic Networks & Intuitive relationship structure & Taxonomies, concept maps, explainability & Semantics can be less strict \\
        \bottomrule
    \end{tabularx}
\end{frame}

\begin{frame}{Ontologies: Basic Idea}
    Ontology = explicit conceptual model of a domain:
    \begin{itemize}
        \item \textbf{Classes:} Patient, Disease, Medication
        \item \textbf{Relations:} hasSymptom, treatedBy, contraindicatedWith
        \item \textbf{Instances:} specific entities (e.g., Patient\_102)
    \end{itemize}
    \vspace{0.2cm}
    Purpose:
    \begin{itemize}
        \item Shared vocabulary across teams/systems.
        \item Better interoperability and data integration.
    \end{itemize}
\end{frame}

\section{Putting It Together}

\begin{frame}{End-to-End KR Pipeline (Practical)}
    \begin{enumerate}
        \item Define domain concepts and relations.
        \item Choose representation (logic/rules/net/ontology blend).
        \item Encode facts and rules.
        \item Select inference strategy.
        \item Validate outputs with domain experts.
        \item Iterate as domain evolves.
    \end{enumerate}
\end{frame}

\begin{frame}{Common Pitfalls in KR Projects}
    \begin{itemize}
        \item \textbf{Over-modeling:} too much detail too early.
        \item \textbf{Ambiguous predicates:} unclear symbol meaning.
        \item \textbf{Hidden assumptions:} not encoded explicitly.
        \item \textbf{No governance:} rules drift and become inconsistent.
    \end{itemize}
    \vspace{0.2cm}
    Mitigation: naming conventions, review cycles, test queries, and versioning.
\end{frame}

\section{Assessment and Wrap-Up}

\begin{frame}{Assessment Q1 (Conceptual)}
    \textbf{Question:} Why is "representation" as important as the inference algorithm?

    \vspace{0.2cm}
    \textbf{Expected discussion points:}
    \begin{itemize}
        \item It determines what can be expressed.
        \item It constrains what can be inferred efficiently.
        \item Poor representation can make good inference unusable.
    \end{itemize}
\end{frame}

\begin{frame}{Assessment Q2 (Logic Modeling)}
    Convert to propositional statements and infer:
    \begin{itemize}
        \item If server overload then alert admin.
        \item If alert admin then open incident ticket.
        \item Server overload occurred.
    \end{itemize}
    \vspace{0.2cm}
    \textbf{Ask students:} What follows logically? What does not necessarily follow?
\end{frame}

\begin{frame}{Assessment Q3 (FOL Modeling)}
    \textbf{Task:} Write FOL for:
    \begin{itemize}
        \item Every student in this course has an ID.
        \item Riya is a student in this course.
    \end{itemize}
    \vspace{0.2cm}
    Then ask: what can be concluded about Riya?
\end{frame}

\begin{frame}{Assessment Q4 (Inference Strategy)}
    A doctor asks: \textit{"Does patient P likely have condition C?"}

    \vspace{0.2cm}
    \textbf{Prompt:}
    \begin{itemize}
        \item Which is more natural first: forward chaining or backward chaining?
        \item Why?
    \end{itemize}
\end{frame}

\begin{frame}{Assessment Q5 (Representation Choice)}
    For a university knowledge system (courses, prerequisites, instructors, departments):
    \begin{itemize}
        \item What would you store as rules?
        \item What would you store as ontology classes/relations?
        \item Where would semantic network visualization help stakeholders?
    \end{itemize}
\end{frame}

\begin{frame}{Key Takeaways}
    \begin{enumerate}
        \item Knowledge representation and inference are inseparable in AI design.
        \item Propositional logic gives clean foundations; FOL adds object-level expressiveness.
        \item Chaining and resolution provide complementary inference perspectives.
        \item Rules, semantic nets, and ontologies serve different engineering needs.
        \item Real systems often combine these representations, not just one.
    \end{enumerate}
\end{frame}

\begin{frame}[standout]
    Questions?
\end{frame}

\end{document}
